\section{Fixed-coefficient cyclotomic transport}

\begin{proposition}[AC-K33-cyclotomic: one denominator bound for the fixed quotient]
Fix the quotient, tame principal level and coefficient lattices of
AC-K33.  Its base change to every \(\mathbf Q(\mu_m)\) commutes
with restriction and corestriction, and one fixed scalar makes all
these coefficient maps integral.  Here \(m\) ranges over positive
integers prime to \(pMN\).  The preordinary geometric transport can
be performed on the actual full-tame curves: there is a map
\[
 s_{m;M,N}:Y(Mm,Nm)\longrightarrow
       Y(M,N)\times\mu_m^\circ
\]
which is compatible with internal \(p\)-level transitions and
auxiliary cyclotomic trace.  These assertions concern the fixed
Artin coefficient over cyclotomic extensions.  They do not construct
the varying CM coefficients \(\theta_{\chi\psi}\) over general
\(K\)-ray extensions or remove either factor of a split-prime
Rankin Euler polynomial.
\end{proposition}

\begin{proof}
For \((E,P,Q)\in Y(Mm,Nm)\), let
\(C=\langle MP\rangle\), a cyclic subgroup of order \(m\), and let
\(\phi:E\to E'=E/C\).  Define
\[
 s_{m;M,N}(E,P,Q)=
 \left((E',\phi(P),\phi(mQ)),\
       \langle MP,NQ\rangle_{E[m]}\right).
\]
The output points have exact orders \(M\) and \(N\).
If a linear combination of them vanishes, the corresponding
combination of \(P\) and \(mQ\) lies in \(\langle MP\rangle\);
independence in the original level structure forces its two
coefficients to be divisible by \(M\) and \(N\).
The displayed Weil pairing is primitive.  Thus this is a morphism
of the stated moduli problems over \(\mathbf Q\), with all
cyclotomic components retained.

The complex uniformization sends \(\tau\) to \(\tau/m\).
Conjugation by \(\operatorname{diag}(1,m)\) sends
\(\Gamma(Mm,Nm)\) into \(\Gamma(M,N)\).
This also gives the finite covering description.  On the
Tate division torsor the map is multiplication by \(m\),
followed by the isogeny \(\phi\), as in the coefficient
normalization of KLZ Definition~6.1.1.
On \(\Lambda^{[j]}=\Lambda\otimes\operatorname{TSym}^j\mathcal H\),
the multiplication \([m]\) acts only on the torsor/measure
factor; the algebraic \(\operatorname{TSym}^j\mathcal H\)
factor is transported by \(\operatorname{TSym}^j(\phi)\).
This is the convention of KLZ Notation~5.2.1 and
Definition~6.1.1.  In particular no extra scalar \(m^j\)
is inserted into the Clebsch--Gordan norm relation.
Since \(p\nmid m\), these fibre maps are integral
isomorphisms.
For the internal \(p\)-level map,
\(\phi(mpQ)=p\phi(mQ)\), so compatibility is exact.

For \(m\to mq\), first quotient by
\(\langle MmP\rangle\) of order \(q\), then by the
image of \(\langle MP\rangle\) of order \(m\).
The total kernel is \(\langle MP\rangle\) of order \(mq\).
The second marked point is multiplied by \(mq\).
Hence the two composite curve and coefficient maps agree;
the Weil-pairing component changes by
\(\zeta_{mq}\mapsto\zeta_{mq}^{q}\).
Functoriality of finite pushforward gives the cyclotomic
trace compatibility.  Although the degrees of the finite
coverings may have \(p\)-factors, traces are integral:
no inverse degree is used.  On the two-factor product
one uses the fibre product over the common
\(\mu_m^\circ\)-component.  The Rankin--Iwasawa diagonal
and its unipotent translate have equal Weil-pairing
components, so their classes lie in this open-and-closed
part of the product.

Finally, a fixed bounded \(G_{\mathbf Q}\)-map from the
compact coefficient module has bounded image in its
finite-dimensional target.  Its image spans a fixed
\(G_{\mathbf Q}\)-stable \(\mathcal O\)-lattice.
Alternatively, relative to any prechosen stable lattice,
one fixed power of \(p\) makes the map integral.
After any field extension the coefficient map is the
same map on continuous Galois cochains.  The finite sum
defining corestriction commutes with it term by term.
Tensoring with a fixed first-factor quotient adds only
its one fixed denominator.  All geometric operations
above are integral, so the bound is independent of \(m\).

None of these operations replaces the rational norm
polynomial at a split auxiliary prime by one of its
two \(K\)-prime factors.  In particular no compatibility
for the varying family \(\chi\psi\), and no nonzero
projected Rankin--Iwasawa class, follows from boundedness.
\end{proof}

\begin{remark}[Printed normalization interface]
The coefficient definition in KLZv5 \S6.1.1 uses
multiplication by the auxiliary modulus followed by
the quotient isogeny.  The construction above follows
that coefficient normalization and independently checks
the orders of the two marked points and the adelic
matrix.  The preceding v5 displayed moduli sentence
omits that multiplication on its second marked point;
no inference in this proposition relies on that
sentence without the order check.  This source-interface
note does not affect AC-K33's distinct use of
Theorem~4.5.1.
\end{remark}
